\chapter{Memilih Model dengan Statistik dan Bayes}
\label{ch:stat}

Pada musim dingin 2025/26 saya mengikuti kuliah Machine Learning: Basic Techniques dari Sepp Hochreiter, kuliah yang
namanya ``dasar'' tetapi dalam perjalanan saya datang setelah sebagian besar kuliah lanjut. Naskahnya dibuka dengan satu klaim yang terus diulang sampai halaman terakhir: inti machine learning
adalah \emph{model selection}. Memilih arsitektur, memilih fungsi loss, memilih regularisasi, bahkan memilih cara meringkas
data sebelum memodelkannya, semuanya adalah memilih model, dan setiap pilihan itu menyimpan asumsi tentang dunia. Bab ini
mengikuti alur naskah itu. Ia melengkapi \cref{ch:data} dengan sisi statistik yang lebih klasik, yang sering terlewat
ketika orang langsung melompat ke jaringan saraf.

\section{Deduktif dan induktif}

Naskah kuliah membedakan dua cara membangun model. Model deduktif diturunkan dari prinsip pertama: hukum Newton, persamaan
Navier--Stokes, aturan yang ditulis ahli. Model induktif dibangun dari contoh: kita memilih sebuah keluarga model, lalu
data menentukan anggota keluarga mana yang dipakai. Machine learning adalah pembangunan model induktif. Bab-bab tentang
simulasi di bagian akhir buku ini pada dasarnya adalah cerita tentang mempertemukan kedua cara itu.

Keluarga model sendiri bisa parametrik atau nonparametrik. Model parametrik punya jumlah parameter tetap, berapa pun
datanya, seperti regresi linear atau jaringan saraf berukuran tertentu. Model nonparametrik tumbuh bersama data, seperti
$k$-nearest neighbour yang menyimpan semua contoh atau penaksir kerapatan kernel di bawah. Model parametrik kuat ketika
asumsinya benar dan berbahaya ketika salah. Model nonparametrik lebih lentur, dengan harga butuh data lebih banyak.

\section{Meringkas data sebelum memodelkan}

Naskah kuliah menghabiskan satu bab penuh untuk hal yang tampak sepele: meringkas dan menggambar data. Contohnya data bunga
iris milik Anderson dan data microarray dari beberapa jaringan tubuh. Alasannya sederhana. Ringkasan yang salah memberi
gambaran yang salah, dan model yang dibangun di atas gambaran yang salah tidak akan menjadi benar karena arsitekturnya
canggih.

\subsection{Pusat dan sebaran}

Rata-rata, median, dan modus sama-sama mengukur pusat, tetapi bereaksi berbeda terhadap pencilan. Lima gaji bulanan, dalam
juta rupiah, $3$, $4$, $4$, $5$, dan $100$, punya rata-rata $23{,}2$ dan median $4$. Satu angka ekstrem menarik rata-rata jauh
dari tempat sebagian besar data berada, sedangkan median nyaris tidak terganggu. Median disebut \istilah{robust}: setengah
data harus dirusak sebelum median bisa digeser sembarang jauh, sedangkan untuk rata-rata, satu titik saja sudah cukup.

Untuk sebaran, varians sampel bisa dihitung dengan pembagi $N$ atau $N-1$. Pembagi $N$ memberi penaksir yang bias. Alasannya,
simpangan dihitung terhadap rata-rata sampel, bukan rata-rata sejati, dan rata-rata sampel selalu terletak tepat di tengah
data itu sendiri, sehingga jumlah kuadrat simpangannya selalu sedikit lebih kecil dari seharusnya. Menghitung nilai
harapannya memberi $\E\big[\tfrac1N\sum(x_i-\bar x)^2\big]=\tfrac{N-1}N\sigma^2$, dan membagi dengan $N-1$ menghapus bias itu.
Tetapi naskah kuliah menunjukkan sesuatu yang lebih halus. Untuk data Gauss, penaksir yang bias justru punya galat kuadrat
rata-rata yang lebih kecil, karena variansnya sendiri lebih kecil, dan pengurangan varians itu lebih besar dari tambahan
bias kuadratnya. Tak bias tidak sama dengan terbaik. Inilah tawar-menawar bias dan varians dari \cref{ch:data}, dalam
bentuk yang paling kecil.

Pilihan antara rata-rata dan median juga bergantung pada distribusinya. Kalau data berdistribusi Gauss, varians rata-rata
sampel adalah $\sigma^2/N$, sedangkan varians median sampel kira-kira $\tfrac\pi2\,\sigma^2/N$, sekitar 57 persen lebih besar.
Untuk data Gauss, rata-rata adalah penaksir pusat yang lebih baik. Kalau data berdistribusi Laplace, yang ekornya lebih tebal,
keadaannya terbalik: varians median sampel hanya $\sigma^2/(2N)$, separuh varians rata-rata. Tidak ada penaksir yang terbaik
untuk semua keadaan. Yang ada hanya penaksir yang cocok untuk asumsi tertentu.

\subsection{Menggambar data}

\istilah{Boxplot} meringkas satu variabel dengan lima angka. Kotaknya membentang dari kuartil pertama ke kuartil ketiga, garis di
dalamnya adalah median, dan kumis memanjang sampai titik data terjauh yang masih berada dalam $1{,}5$ kali jarak antarkuartil
dari kotak. Titik di luar itu digambar sendiri sebagai calon pencilan. Boxplot per kelas pada data iris langsung
memperlihatkan bahwa panjang mahkota bunga memisahkan spesies jauh lebih baik dari panjang kelopaknya, sebelum satu model
pun dilatih.

Histogram bergantung pada lebar dan letak kotak-kotaknya, sehingga bentuknya bisa berubah hanya karena kotaknya digeser.
Penaksir kerapatan kernel menghaluskannya dengan meletakkan sebuah kernel, biasanya Gauss, di setiap titik data dan
menjumlahkannya,
\begin{equation}
\hat p(x)=\frac1{Nh}\sum_{i=1}^NK\Big(\frac{x-x_i}h\Big).
\end{equation}
Lebar pita $h$ memainkan peran yang sudah kita kenal. Pita yang terlalu sempit membuat setiap titik data menjadi puncak
sendiri, overfitting. Pita yang terlalu lebar meratakan dua puncak menjadi satu, underfitting. Violin plot menggabungkan
boxplot dengan kerapatan kernel yang dicerminkan, sehingga bentuk distribusi dan ringkasan lima angkanya terlihat sekaligus.

\subsection{Korelasi dan bahaya angka tunggal}

Koefisien korelasi Pearson mengukur seberapa linear hubungan dua variabel. Untuk menguji apakah korelasi sampel $r$ dari $N$
pasangan berbeda nyata dari nol, dengan asumsi kedua variabel berdistribusi normal, dipakai statistik
\begin{equation}
t=r\sqrt{\frac{N-2}{1-r^2}},
\end{equation}
yang mengikuti distribusi Student-$t$ dengan $N-2$ derajat kebebasan. Naskah kuliah memakai kuartet Anscombe
\citep{anscombe1973}, empat himpunan data berisi sebelas titik yang punya rata-rata, varians, korelasi, dan garis regresi yang
hampir sama persis. Untuk himpunan pertama, $r\approx0{,}816$ dan $t=0{,}816\sqrt{9/0{,}334}\approx4{,}24$, dengan nilai $p$ satu
sisi sekitar $0{,}001$. Tetapi ketika digambar, keempatnya sama sekali berbeda: satu awan titik yang wajar, satu kurva
lengkung yang sempurna, satu garis lurus dengan satu pencilan, dan satu kolom titik dengan satu titik jauh yang sendirian
menciptakan seluruh korelasinya. Angka ringkasan yang sama bisa berasal dari dunia yang jauh berbeda. Gambarlah data
sebelum mempercayai angkanya.

\section{Loss sebagai asumsi noise}

Di \cref{ch:data} kita menurunkan galat kuadrat dari asumsi noise Gauss. Naskah kuliah melanjutkannya ke noise lain. Kalau
noise berdistribusi Laplace, $p(\epsilon)\propto e^{-\beta|\epsilon|}$, log likelihood negatifnya sebanding dengan
$\sum|y_i-\hat y_i|$, galat mutlak. Secara umum, noise dengan kerapatan $\propto e^{-\beta|\epsilon|^r}$ menghasilkan galat
Minkowski $\sum|y_i-\hat y_i|^r$. Pilihan $r$ adalah pilihan seberapa tebal ekor noise yang kita percayai.

Hubungan antara loss dan penaksir pusat bisa diturunkan dalam beberapa baris. Cari satu konstanta $c$ yang mewakili data
$y_1,\dots,y_N$. Untuk galat kuadrat, turunan $\sum(y_i-c)^2$ terhadap $c$ adalah $-2\sum(y_i-c)$, yang nol tepat ketika $c$
adalah rata-rata. Untuk galat mutlak, turunan $\sum|y_i-c|$ adalah banyaknya titik di bawah $c$ dikurangi banyaknya titik di
atas $c$, karena setiap titik menyumbang $+1$ atau $-1$. Turunan itu nol ketika titik di bawah dan di atas sama banyak, yaitu
ketika $c$ adalah median. Galat kuadrat menghasilkan rata-rata, galat mutlak menghasilkan median, dan kita sudah melihat
bahwa median tahan terhadap pencilan. Jadi memilih loss L1 untuk regresi berarti memilih model yang tidak mudah ditarik oleh
beberapa pengukuran yang rusak. Untuk data gaji di atas, model L2 akan meramal $23{,}2$, model L1 meramal $4$.

\section{Batas bawah galat: klasifikasi Bayes}

Seberapa baik sebuah pengklasifikasi bisa bekerja? Kalau distribusi data diketahui, jawabannya bisa dihitung. Pengklasifikasi
yang memilih kelas dengan peluang posterior tertinggi, $\argmax_yp(y\mid\vx)$, mencapai risiko serendah mungkin, dan risiko itu
disebut galat Bayes. Tidak ada model, berapa pun besarnya, yang bisa lebih baik dari ini, karena di daerah tempat kedua kelas
tumpang tindih, informasi untuk membedakannya memang tidak ada.

Naskah kuliah menghitungnya untuk dua kelas yang masing-masing berdistribusi Gauss. Log rasio kedua posterior adalah
\begin{equation}
\ln\frac{p(y=1\mid\vx)}{p(y=-1\mid\vx)}=
-\tfrac12(\vx-\bm\mu_1)^\top\bm\Sigma_1^{-1}(\vx-\bm\mu_1)+\tfrac12(\vx-\bm\mu_{-1})^\top\bm\Sigma_{-1}^{-1}(\vx-\bm\mu_{-1})
+\text{konstanta}.
\end{equation}
Kalau kedua kovarians sama, suku kuadrat $\vx^\top\bm\Sigma^{-1}\vx$ dari kedua kelas saling menghapus, dan yang tersisa linear
dalam $\vx$, dengan arah $\bm\Sigma^{-1}(\bm\mu_1-\bm\mu_{-1})$, arah yang sama dengan analisis diskriminan Fisher di
\cref{ch:unsup}. Batas keputusannya adalah bidang datar. Kalau kovariansnya berbeda, suku kuadrat tidak terhapus, dan batasnya
menjadi permukaan kuadrat, misalnya hiperbola atau elips, seperti yang digambarkan naskah kuliah. Model yang dipilih untuk
batas keputusan, linear atau kuadrat, adalah asumsi tentang kovarians kelas.

Dalam satu dimensi, galat Bayes bisa dihitung dengan tangan (\cref{fig:bayeserror}). Dua kelas dengan peluang sama,
$\N(0,1)$ dan $\N(2,1)$, punya batas keputusan di tengah, $x=1$. Kesalahan terjadi ketika contoh dari kelas pertama jatuh di
atas 1 atau contoh dari kelas kedua jatuh di bawah 1. Masing-masing berpeluang $\Phi(-1)\approx0{,}159$, dengan $\Phi$ fungsi
distribusi kumulatif normal baku, sehingga galat Bayes sekitar 16 persen. Kalau jarak kedua rata-rata dua kali lipat, galatnya
turun menjadi $\Phi(-2)\approx2{,}3$ persen. Pengklasifikasi yang dilaporkan mencapai akurasi 95 persen pada data dengan galat
Bayes 16 persen pasti sedang mengukur sesuatu yang salah, misalnya data uji yang bocor.

\begin{figure}[ht]
\centering
\begin{tikzpicture}
\begin{axis}[width=0.66\linewidth,height=4.6cm,axis lines=left,xmin=-3,xmax=5,ymin=0,ymax=0.45,
  xlabel={\small $x$},ylabel={\small kerapatan},xtick={0,1,2},ytick=\empty]
  \addplot[name path=A,biru,thick,domain=-3:5,samples=120] {exp(-x^2/2)/sqrt(2*pi)};
  \addplot[name path=B,oranye,thick,domain=-3:5,samples=120] {exp(-(x-2)^2/2)/sqrt(2*pi)};
  \addplot[name path=Z,draw=none,domain=-3:5] {0};
  \addplot[biru!30] fill between[of=A and Z,soft clip={domain=1:5}];
  \addplot[oranye!30] fill between[of=B and Z,soft clip={domain=-3:1}];
  \draw[dashed,abu] (axis cs:1,0) -- (axis cs:1,0.43);
  \node[label] at (axis cs:1,0.43) [right] {batas keputusan};
\end{axis}
\end{tikzpicture}
\caption{Galat Bayes untuk dua kelas Gauss dengan rata-rata 0 dan 2. Daerah yang diarsir adalah contoh yang pasti salah
diklasifikasikan oleh pengklasifikasi terbaik sekalipun.}
\label{fig:bayeserror}
\end{figure}

\section{Memilih model dengan cara Bayes}

Bab terakhir naskah kuliah membawa pemilihan model ke kerangka Bayes sepenuhnya. Parameter $\vw$ diberi prior, data memberi
likelihood, dan teorema Bayes memberi posterior,
\begin{equation}
p(\vw\mid\mathcal D)=\frac{p(\mathcal D\mid\vw)\,p(\vw)}{p(\mathcal D)},\qquad
p(\mathcal D)=\int p(\mathcal D\mid\vw)\,p(\vw)\dd\vw .
\end{equation}
Penyebutnya, peluang data setelah semua parameter yang mungkin dirata-ratakan, disebut \istilah{evidence}. Ia tidak berperan
dalam mencari $\vw$, tetapi menjadi tokoh utama ketika yang dipilih adalah modelnya sendiri.

\subsection{MAP dan error bar}

Memaksimalkan posterior memberi penaksir MAP, yang untuk noise Gauss dan prior Gauss tidak lain adalah galat kuadrat
ditambah weight decay (\cref{ch:data}). Tetapi posterior memberi lebih dari satu titik. Ia juga memberi ketidakpastian, dan
ketidakpastian itu bisa diteruskan ke prediksi. Kalau keluaran model dilinearkan di sekitar $\vw_{\mathrm{map}}$ dan posterior
didekati dengan Gauss yang kovariansnya kebalikan Hessian $\bm H$, naskah kuliah menurunkan varians prediksi
\begin{equation}
\sigma^2(\vx)=\frac1\beta+\bm g^\top\bm H^{-1}\bm g,
\end{equation}
dengan $1/\beta$ varians noise dan $\bm g$ gradien keluaran terhadap parameter di $\vx$. Suku pertama adalah noise yang tidak
bisa dihindari. Suku kedua adalah ketidakpastian tentang parameter, dan ia membesar di tempat model peka terhadap parameter
yang kurang ditentukan data.

Contoh dengan satu parameter membuatnya jelas. Model $y=wx$ dengan prior $w\sim\N(0,1/\alpha)$, $\alpha=1$, noise dengan
$\beta=4$ (simpangan baku $0{,}5$), dan dua titik data di $x=1$ dan $x=2$. Posterior $w$ adalah Gauss dengan presisi
$\alpha+\beta\sum x_i^2=1+4\cdot5=21$, sehingga variansnya $1/21\approx0{,}048$. Di sini $g=x$, sehingga varians prediksi
$\sigma^2(x)=0{,}25+x^2/21$. Di $x=1$, simpangan baku prediksi $\sqrt{0{,}298}\approx0{,}55$, hampir hanya noise. Di $x=3$, di luar
data, nilainya $\sqrt{0{,}679}\approx0{,}82$. Pita ketidakpastian melebar ketika model diminta berbicara jauh dari tempat ia
pernah melihat data (\cref{fig:errorbar}), sifat yang persis kita inginkan dari model mana pun yang dipakai untuk mengambil
keputusan.

\begin{figure}[ht]
\centering
\begin{tikzpicture}
\begin{axis}[width=0.64\linewidth,height=4.8cm,axis lines=left,xmin=-1,xmax=5,ymin=-2.5,ymax=7,
  xlabel={\small $x$},ylabel={\small $y$},xtick={0,1,2,3,4},ytick={0,2,4,6}]
  \addplot[name path=U,draw=none,domain=-1:5,samples=60] {0.933*x+2*sqrt(0.25+x^2/21)};
  \addplot[name path=L,draw=none,domain=-1:5,samples=60] {0.933*x-2*sqrt(0.25+x^2/21)};
  \addplot[biru!18] fill between[of=U and L];
  \addplot[biru,thick,domain=-1:5] {0.933*x};
  \addplot[only marks,mark=*,oranye] coordinates {(1,1.1) (2,1.9)};
\end{axis}
\end{tikzpicture}
\caption{Prediksi model $y=wx$ dengan pita dua simpangan baku dari kerangka Bayes. Pita tersempit di sekitar data dan
melebar ke arah yang jauh dari data, karena ketidakpastian parameter dikalikan $x$.}
\label{fig:errorbar}
\end{figure}

\subsection{Evidence dan pisau Occam}

Hyperparameter seperti $\alpha$ dan $\beta$ juga bisa dipilih dengan Bayes: pilih nilai yang memaksimalkan evidence
$p(\mathcal D\mid\alpha,\beta)$. Kerangka ini, yang dikembangkan \citet{mackay1992} dan dibahas rinci di naskah kuliah, disebut
\istilah{evidence framework}. Untuk model linear dengan noise dan prior Gauss, log evidence punya bentuk tertutup yang berisi
galat latih, suku regularisasi, dan $-\tfrac12\ln|\bm H|$, sehingga tidak perlu data validasi terpisah.

Evidence juga membandingkan model yang berbeda, dan cara kerjanya mengandung pisau Occam yang otomatis. Kalau posterior
memuncak tajam, evidence kira-kira sama dengan likelihood terbaik dikali rasio lebar posterior terhadap lebar prior:
\begin{equation}
p(\mathcal D\mid\mathcal M)\approx p(\mathcal D\mid\vw_{\mathrm{map}},\mathcal M)\cdot\frac{\Delta w_{\mathrm{posterior}}}{\Delta w_{\mathrm{prior}}}.
\end{equation}
Faktor kedua, faktor Occam, selalu lebih kecil dari satu, dan semakin kecil kalau model punya banyak parameter yang harus
disempitkan dari prior yang lebar. Misalkan satu parameter punya prior selebar 10 dan posterior selebar $0{,}1$. Faktor
Occam-nya $0{,}01$, setara dengan menghukum log evidence sebesar $4{,}6$. Model yang lebih rumit harus memperbaiki log
likelihood terbaiknya lebih dari $4{,}6$ untuk setiap parameter seperti itu sebelum ia dipilih. Model yang terlalu sederhana
kalah karena likelihood-nya buruk, model yang terlalu rumit kalah karena faktor Occam-nya, dan evidence memilih di
tengah-tengah tanpa pernah menyentuh data uji.

Kalau posterior tidak berbentuk Gauss, integralnya didekati dengan sampling, misalnya dengan MCMC dari \cref{ch:prob}. Naskah
kuliah juga mencatat satu jebakan untuk jaringan saraf: posteriornya punya banyak puncak yang setara, karena menukar dua
neuron tersembunyi beserta bobotnya tidak mengubah fungsi jaringan, sehingga perkiraan evidence perlu dikoreksi dengan
faktor jumlah permutasi.

\section{Di luar kelas}

Gagasan-gagasan klasik di bab ini muncul kembali di banyak tempat dalam buku. Loss yang tahan pencilan, seperti galat mutlak,
berguna setiap kali sebagian pengukuran bisa rusak, dari sensor yang dikenakan sampai citra medis. Ringkasan robust seperti median dan normalisasi kuantil dipakai pada data microarray
(\cref{ch:genomik}). Error bar dari posterior adalah dasar Gaussian process (\cref{ch:mlat}), dan pita yang melebar jauh dari data
adalah sifat yang kita harapkan dari model pengganti simulasi, yang justru paling berbahaya ketika diam-diam berekstrapolasi.
Dan kuartet Anscombe adalah pengingat yang pantas ditempel di dinding setiap laboratorium: gambarlah datanya.

\sumberbab{Bab ini mengikuti naskah kuliah Machine Learning: Basic Techniques (pemilihan model, ringkasan dan visualisasi
data, teori peluang, risiko minimum untuk klasifikasi Gauss, loss dan asumsi noise, teknik Bayes untuk prior, error bar,
evidence, dan perbandingan model, serta lampiran tentang rata-rata dan median distribusi simetris). Bacaan pendamping:
\citet{bishop2006} dan \citet{mackay1992}.}
